\subsection{Hyperjump за \( \mathcal{O}(N^2) \)}

\subsubsection{Описание решения}

Переформулировав условие задачи, для данного массива нужно найти отрезок с максимальной суммой его элементов.

Учитывая ограничение \( 0\leq N\leq 6\cdot 10^4 \), можно размышлять в сторону алгоритмов квадратичной сложности. Это означает, что мы можем перебрать всевозможные начала и концы отрезков. Но этого пока недостаточно для нахождения суммы элементов на конкретном отрезке.

Для решения такой задачи необходимо предварительно кешировать префиксные суммы элементов массива за линейное время. Это позволит оптимизировать расчёт суммы элементов отрезка с линейного времени до константного.

Остаётся фиксировать текущий максимум суммы на каждой итерации, чтобы получить глобальный максимум.

\subsubsection{Асимптотическая сложность}

\textit{Временная сложность решения} --- квадратичная \( \mathcal{O}(N^2) \). В начале работы алгоритма мы за линейное время создаём массив префиксных сумм и заполняем его. Затем для каждой неупорядоченной пары элементов массива выполняем константное количество операций.

\textit{Пространственная сложность решения} --- линейная \( \mathcal{O}(N) \). В решении используются два массива длины \( N \) и \( N+1 \). При помощи них производятся вычисления, требующие константное количество памяти.

\subsubsection{Листинг кода}

\begin{lstlisting}
#include <iostream>
#include <vector>

int main() {
  int N;
  std::cin >> N;
  std::vector<int> arr;
  std::vector<long long> prefix;
  prefix.push_back(0);
  for (int i = 0; i < N; i++) {
    arr.push_back(0);
    prefix.push_back(0);
  }
  for (int i = 0; i < N; i++) {
    std::cin >> arr[i];
    prefix[i + 1] = prefix[i] + arr[i];
  }
  long long ans = 0;
  for (int i = 1; i <= N; i++) {
    for (int j = i; j <= N; j++) {
      long long curr = prefix[j] - prefix[i - 1];
      if (curr > ans) {
        ans = curr;
      }
    }
  }
  std::cout << ans;
  return 0;
}
\end{lstlisting}